\section{Introduction}
\label{sec:intro}

\IEEEPARstart{M}{onocular} metric depth models \cite{est_depthpro,est_metric3dv2,est_unidepth,est_dav2} are
accurate enough to place a surface from one image, and one depth map per frame can supervise Gaussian
splatting \cite{gs_3dgs,gs_2dgs,gs_dnsplatter} of an aerial survey without matching a single feature. The
map that results looks complete, and it says nothing about what it knows. Opacity is not confidence: a
surfel on a blank roof is drawn with the same certainty as one on a textured wall.

The danger is an error that repeats. On a given surface a depth model tends to make the same mistake in
every frame, so the depth maps agree with each other and averaging them removes only the part of the error
that differs between frames. If a share $\varrho$ of the error variance is common to the $n$ frames that
see a surface, the variance of their fused estimate is $\varrho+(1-\varrho)/n$ times that of one frame,
whereas the fused information claims $1/n$: the stated variance is too small by $1+(n-1)\varrho$, which
grows without bound with the number of frames. Tests that compare each depth map with the others
\cite{stat_teunissen_dia,flt_raim_joerger,fus_colmap_mvs} are blind to the common part by construction.

The pixels are the only independent witness, and they speak about geometry only through a coupling: an
image constrains a surface if moving the surface would change the image, which needs texture contrast and
views from different directions. The strength of that coupling is a stiffness that can be computed only
where geometry and appearance meet, in the renderer. For a planar surfel \cite{gs_2dgs} the two sources
meet exactly: to first order the derivative of the rendered image with respect to the surfel's placement
factors through the derivative of its depth, so the depth prior and the pixels are quadratic forms on the
same three coordinates, an offset along the normal and two tilts (Fig.~\ref{fig:teaser}). Their
disagreement is then a statistic with units and a null distribution.

Stiffness computed from rendering residuals is not honest as it stands, because residuals are model misfit
and neighbouring pixels are not independent; uncertainty fields derived from the rendering Hessian
\cite{unc_fisherrf,unc_pup3dgs,unc_ougs} are used to rank views or prune primitives, where only the order
matters. We calibrate it on the scene, without labels, from frame replicates and from displacements
injected through the renderer, and we cap the information of the prior at what its repeated error allows.
The disagreement is read as a score test at the surface the depth prior placed, with colour projected from
other frames, so that a wrong surface cannot be hidden by refitting its appearance.

We contribute:
\begin{itemize}
  \item per-source information records on planar surfels, built on the factorisation of the image Jacobian
        through the depth Jacobian; records add over frames and tiles (Section~\ref{sec:method:records});
  \item the overconfidence law $\varrho+(1-\varrho)/n$ for depth-only fusion and a score test at the
        depth-placed surface that detects repeated depth faults, both compared with a depth-consensus gate
        (Sections~\ref{sec:method:honest}--\ref{sec:method:test});
  \item a label-free calibration of the pixel information, scored by a proper rule, the surface log-score
        (Sections~\ref{sec:method:honest} and~\ref{sec:setup}).
\end{itemize}
We do not claim to be more accurate than matching-based pipelines. The pixel coupling has a short capture
range, which the depth prior supplies, and where the pixels have no stiffness the map reports that the
surface is unsupported.
\TODO{Replace the list by result-bearing contributions once the experiments of Section~\ref{sec:results}
are run.}

\begin{figure*}[!t]
\centering
\input{\lfigdir/fig_teaser.tex}
\caption{A depth-only map is confidently wrong where the depth model repeats its error (schematic, surfels of
2~m). (a) Every frame's depth places the larger roof 2.5~m below the true roof (dashed), so the $n=12$ depth
maps agree with each other. (b) The depth-only map, coloured by its standardised error $z$: its stated
$2\sigma$ slab (black) is thin and lies wholly below the true roof. (c) The per-source map, coloured by the
energy $E_m$ between the depth prior and the pixels: the energy is high only on the faulty roof, the depth
record is left out there, and the honest $2\sigma$ slab (red) contains the true roof on both buildings.}
\label{fig:teaser}
\end{figure*}
